\documentclass[11pt]{article}
\usepackage{amsmath,amssymb,amsthm}
\usepackage[margin=1in]{geometry}

\title{V18 P6-T02 Finite Toy Response v2 Model-or-Obstruction}
\author{Candidate Constructor packet RT-20260708-015}
\date{2026-07-08}

\newtheorem{definition}{Definition}
\newtheorem{lemma}{Lemma}
\newtheorem{proposition}{Proposition}

\begin{document}
\maketitle

\section*{Control metadata}

\begin{verbatim}
plan_id: recommendations_implementation_plan_continue_task-v18
plan_phase_id: P6
plan_task_id: P6-T02
task_type: finite_toy_response_v2_model_or_obstruction
task_id: RT-20260708-015
job_id: AJ-RT-20260708-015-001
role: candidate-constructor@0.2.0
model_or_obstruction_result: positive_toy_model_constructed
positive_toy_model: true
precise_obstruction: false
obstruction_id: none
target_derivation_milestone: finite_toy_metric_response
milestone_burden: Construct a non-tag finite toy response model or record a precise obstruction.
claim_status: draft/control
promotion_allowed: false
no_target_metric_import: true
not_g_eff: true
not_matter_coupling: true
not_einstein_equations: true
not_benchmark_promotion: true
next_plan_task_id: P6-T03
\end{verbatim}

\section{Question}

P6-T02 asks whether the P6-T01 source specification supports a finite toy
response model without reviving the frozen explicit-tag route. The required
answer is exactly one of:
\[
  \text{positive finite toy model}
  \quad\text{or}\quad
  \text{precise scoped obstruction}.
\]
This packet records the positive branch.

\section{Constructed finite toy response model}

\begin{definition}[Finite source object]
Let
\[
  S_{v2}=\{a,b,c\},\qquad
  A_{v2}=\bigl\{\{a,b\},\{b,c\}\bigr\}.
\]
The incidence relation is
\[
  Inc_{v2}(x,e) \iff x\in e,\quad x\in S_{v2},\ e\in A_{v2}.
\]
The source object is the finite path
\[
  P_{v2}:=(S_{v2},A_{v2},Inc_{v2}).
\]
\end{definition}

\begin{verbatim}
finite_source_set: S_v2 = {a,b,c}
source_relations:
  - A_v2 = {{a,b},{b,c}}
  - Inc_v2(x,e) iff x in e
\end{verbatim}

\begin{definition}[Source-local distance and response]
For \(x,y\in S_{v2}\), define \(d_A(x,y)\) to be the minimum number of
adjacency edges in an \(A_{v2}\)-path from \(x\) to \(y\). Since \(P_{v2}\) is
connected, \(d_A\) is total on \(S_{v2}\times S_{v2}\). Define
\[
  R_{v2}(\{x,y\}) := d_A(x,y),
  \qquad
  D_{v2}(x,y) := R_{v2}(\{x,y\}).
\]
Thus \(R_{v2}\) has values \(0,1,2\), induced only by the source adjacency
relation.
\end{definition}

\begin{verbatim}
induced_response_relation: R_v2({x,y}) = d_A(x,y)
metric_response_analogue: D_v2(x,y) = R_v2({x,y}) as finite graph-distance form
no_target_metric_import: true
not_g_eff: true
\end{verbatim}

\begin{definition}[Orbit/invariant structure]
Let \(Aut(P_{v2})\) be the automorphism group of the finite path. It contains
the identity and the reflection \(\rho(a)=c,\rho(b)=b,\rho(c)=a\). On
unordered pairs with repetition, the orbit classes are
\[
  O_0=\{\{a,a\},\{b,b\},\{c,c\}\},\quad
  O_1=\{\{a,b\},\{b,c\}\},\quad
  O_2=\{\{a,c\}\}.
\]
The response relation factors through these orbit-distance classes:
\[
  \{x,y\}\in O_k \implies R_{v2}(\{x,y\})=k,\qquad k\in\{0,1,2\}.
\]
\end{definition}

\begin{verbatim}
orbit_or_invariant_structure:
  automorphism_group: Aut(P3) = {id, reflection}
  unordered_pair_orbits:
    - O_0 diagonal/source-coincidence class
    - O_1 adjacent-pair class
    - O_2 endpoint-pair class
source_local_token_candidates: orbit-distance classes only
explicit_target_tags_used: false
\end{verbatim}

\section{Orientation, normalization, and token candidates}

The construction uses only source-local candidates:

\begin{itemize}
  \item The orientation candidate is neutral: no directed edge orientation is
  selected. Reflection is an allowed source automorphism, not a broken
  symmetry.
  \item The normalization candidate is the source edge-count unit
  \(d_A(x,y)=1\) for adjacent vertices.
  \item The token candidates are the orbit-distance classes \(O_0,O_1,O_2\).
  Their names are presentation names for invariant classes; they are not
  independent target response tags.
\end{itemize}

\section{Tag independence}

\begin{lemma}[Presentation-label erasure]
Erase the vertex names \(a,b,c\) but retain the isomorphism class of the
three-vertex path relation. Then the orbit-distance response relation is still
defined up to unique graph isomorphism.
\end{lemma}

\begin{proof}
The finite path has one vertex of degree two and two vertices of degree one.
This degree pattern identifies the center and endpoint orbit without using the
presentation labels. The source adjacency relation determines shortest-path
edge count on the unlabeled path, and the pair orbits are exactly the
distance classes \(0,1,2\). Hence \(R_{v2}\) is recoverable from the source
relation structure alone.
\end{proof}

\begin{verbatim}
tag_independence_argument: presentation labels can be erased; the unlabeled
  path relation still determines center/endpoint orbits and the response
  classes 0,1,2.
explicit_tag_route_reintroduced: false
\end{verbatim}

\section{Relabeling invariance}

\begin{proposition}[Relabeling invariance]
For every source automorphism \(\sigma\in Aut(P_{v2})\),
\[
  R_{v2}(\{\sigma x,\sigma y\})=R_{v2}(\{x,y\})
  \quad\text{for all }x,y\in S_{v2}.
\]
\end{proposition}

\begin{proof}
An automorphism of \(P_{v2}\) preserves adjacency and therefore preserves path
length. Hence \(d_A(\sigma x,\sigma y)=d_A(x,y)\). The result follows from
\(R_{v2}(\{x,y\})=d_A(x,y)\).
\end{proof}

\begin{verbatim}
relabeling_invariance_argument: every source automorphism preserves adjacency,
  path length, orbit classes, and therefore R_v2.
\end{verbatim}

\section{Finite-variation probe for the next stress packet}

P6-T03 owns adversarial stress. This packet nevertheless records a bounded
probe target so the next Refuter has a concrete object to attack. For a
single finite edit \(A'\) of \(A_{v2}\) on the same source set, define
\[
  R_{v2}^{A'}(\{x,y\})=
  \begin{cases}
    d_{A'}(x,y), & \text{if }x,y\text{ lie in the same }A'\text{-component},\\
    \bot, & \text{otherwise.}
  \end{cases}
\]
The value \(\bot\) is a fail-closed source-local obstruction marker for a
disconnected pair, not a target metric value and not an empirical detector
readout. The base model remains the connected path model above.

\section{Claim boundary}

This is a finite toy response model only. It does not construct or adopt a
physical metric, matter coupling, detector/readout semantics, stress-energy
semantics, matter action, field equations, benchmark status, or any completed
derivation claim.

\begin{verbatim}
not_matter_coupling: true
detector_readout_status: placeholder_only_nonadopted
source_law_adopted: false
canonical_ontology_edit: false
distance_to_gr_ledger_delta: false
benchmark_status_changed: false
\end{verbatim}

\section{Candidate constructor result}

\begin{verbatim}
candidate_constructor_result:
  result_type: constructed_candidate
  constructed_candidate_path: research_control/tasks/RT-20260708-015/artifacts/finite_toy_response_v2_model_or_obstruction.tex
  formal_objects:
    - P_v2 = (S_v2,A_v2,Inc_v2)
    - Aut(P_v2) orbit-distance classes O_0,O_1,O_2
    - R_v2({x,y}) = d_A(x,y)
    - finite-variation fail-closed probe R_v2^{A'}
  maps:
    - source adjacency relation -> graph distance d_A
    - graph distance -> response relation R_v2
    - source automorphism orbits -> invariant response classes
  proof_obligations:
    - P6-T03 relabeling and tag-removal stress
    - P6-T03 finite-variation fail-closed stress
    - later selector must decide source-model zoo integration without promotion
  no_fog_check: true
  next_required_role: refuter
  claim_boundary_preserved: true
\end{verbatim}

\section{Conclusion}

P6-T02 is satisfied by the positive construction above. The construction is
source-local, finite, relation-induced, relabeling invariant, and independent
of presentation labels. It remains draft/control and must be stress-tested by
P6-T03 before any selector may integrate it into the source-model zoo.

\end{document}
